\notechapter{N-20221223}{Homotopy, Simply Connected Spaces, Semigroups, Groups, and Homomorphisms}

\section{Homotopy intuition}

The first part explains homotopy through paths in \(\mathbb C\).  If \(p,q\in\mathbb C\), many paths from \(p\) to \(q\) can be continuously deformed into one another while the endpoints stay fixed.  But in
\[
  \mathbb C\setminus\{0\},
\]
the missing origin becomes an obstacle.  A path going around the hole one way may not be deformable into a path going the other way without crossing the missing point.

A homotopy between paths \(f_0,f_1:[0,1]\to X\) is a continuous map
\[
  F:[0,1]\times[0,1]\to X
\]
with \(F(s,0)=f_0(s)\), \(F(s,1)=f_1(s)\).  For fixed-endpoint homotopy, the endpoints remain fixed throughout.

\section{Simply connected spaces}

A path-connected space is simply connected if every loop can be contracted to a point.  The plane \(\mathbb C\) is simply connected; the punctured plane \(\mathbb C\setminus\{0\}\) is not.  This is the fundamental intuition behind algebraic topology: holes can be detected by loops that refuse to shrink.

\section{Semigroups, monoids, and groups}

A semigroup is a set with an associative binary operation.  A monoid is a semigroup with an identity element.  A group is a monoid in which every element has an inverse:
\[
  \text{semigroup}\longrightarrow\text{monoid}\longrightarrow\text{group}.
\]
A group \(G\) has associativity, identity, and inverses.  If \(ab=ba\) for all \(a,b\in G\), it is abelian.  Examples include \((\mathbb Z,+)\), \(S_n\), and symmetry groups of geometric figures.

\section{Congruence, quotients, and homomorphisms}

A congruence relation on an algebraic structure is an equivalence relation compatible with the operation.  This lets operations descend to equivalence classes.  For groups, normal subgroups play this role for quotient groups.

A group homomorphism \(f:G\to H\) satisfies
\[
  f(ab)=f(a)f(b).
\]
It automatically gives \(f(e_G)=e_H\) and \(f(a^{-1})=f(a)^{-1}\).  Its kernel
\[
  \ker f=\{g\in G:f(g)=e_H\}
\]
is a normal subgroup, and the image is a subgroup.


\section{Loops and the beginning of the fundamental group}

A loop based at \(x_0\in X\) is a path
\[
  \gamma:[0,1]\to X,\qquad \gamma(0)=\gamma(1)=x_0.
\]
Two loops are considered equivalent if one can be continuously deformed into the other while keeping the basepoint fixed.  The set of such equivalence classes becomes a group under concatenation:
\[
  [\gamma]\cdot[\eta]=[\gamma*\eta].
\]
The identity is the constant loop, and the inverse of a loop is the same path traversed backward.

This explains why group theory appears immediately after homotopy in the note.  The algebra is not an unrelated topic; it is the algebra of loops.

\section{Why the punctured plane has nontrivial loops}

In \(\mathbb C\setminus\{0\}\), the loop
\[
  \gamma(t)=e^{2\pi i t}
\]
winds once around the origin.  It cannot be contracted to a point without crossing the missing origin.  More generally,
\[
  \gamma_n(t)=e^{2\pi i n t}
\]
winds \(n\) times.  The integer \(n\) is the winding number.  This gives the fundamental group
\[
  \pi_1(\mathbb C\setminus\{0\})\cong\mathbb Z.
\]
The plane itself has trivial fundamental group because every loop can be shrunk through the missing-free interior.

\section{Homomorphisms as structure-preserving maps}

The algebraic part of the note introduces homomorphisms because topology will soon produce groups naturally.  A homomorphism does not preserve elements individually; it preserves the operation:
\[
  f(ab)=f(a)f(b).
\]
From this one derives
\[
  f(e_G)=e_H,\qquad f(a^{-1})=f(a)^{-1}.
\]
The kernel records what the map collapses to the identity:
\[
  \ker f=\{g\in G:f(g)=e_H\}.
\]
In topology, continuous maps induce homomorphisms on fundamental groups.  Thus a map of spaces becomes a map of algebraic invariants.

\section{Quotients in topology and algebra}

Both quotient spaces and quotient groups follow the same pattern: impose an equivalence relation and then ask whether the structure descends.  For spaces, open sets in the quotient are defined by preimages under the quotient map.  For groups, normal subgroups are exactly what make multiplication of cosets well-defined.

This shared pattern is one of the first conceptual bridges of algebraic topology.  A loop space quotient by homotopy becomes a group; a group quotient by a normal subgroup becomes another group.  In both cases, equivalence classes are not just bookkeeping but new mathematical objects.

\section{Higher viewpoint: from holes to algebra}

The important transition is:
\[
  \text{hole in a space}\quad\leadsto\quad\text{nontrivial loop}\quad\leadsto\quad\text{group element}.
\]
This is the beginning of algebraic topology.  Instead of trying to draw all possible deformations, we assign algebraic invariants to spaces.  If two spaces have different fundamental groups, they cannot be homeomorphic.

\section{A useful afterword}

The note looks split between topology and algebra, but the split is exactly the point.  Homotopy studies continuous deformation of paths; groups study composable reversible operations.  The fundamental group combines them: homotopy classes of loops form a group, and that group records holes in the space.

\section{Concatenation of loops and why homotopy classes form a group}

Given loops \(\gamma,\eta:[0,1]\to X\) based at \(x_0\), their concatenation is
\[
  (\gamma*\eta)(t)=
  \begin{cases}
  \gamma(2t),&0\le t\le\frac12,\\
  \eta(2t-1),&\frac12\le t\le1.
  \end{cases}
\]
This operation is not strictly associative on the nose, because parametrizations differ.  It becomes associative after passing to homotopy classes.  The identity is the constant loop, and the inverse is
\[
  \overline{\gamma}(t)=\gamma(1-t).
\]
Thus the group structure of \(\pi_1(X,x_0)\) is already a quotient phenomenon: strict path operations become algebraic only after identifying homotopic paths.

\section{The exponential covering and the integer invariant}

The map
\[
  p:\mathbb R\to S^1,\qquad p(t)=e^{2\pi i t}
\]
is a covering map.  A loop in \(S^1\) based at \(1\) lifts uniquely to a path in \(\mathbb R\) starting at \(0\).  The endpoint of the lift is an integer, and that integer is the winding number.

This gives a concrete isomorphism
\[
  \pi_1(S^1)\cong\mathbb Z.
\]
The generator is one counterclockwise turn.  The inverse is one clockwise turn.  Concatenating loops adds winding numbers.

\section{Punctured plane deformation retracts to the circle}

The punctured plane \(\mathbb C\setminus\{0\}\) has the same fundamental group as \(S^1\) because it deformation retracts onto the unit circle:
\[
  H(z,t)=\left((1-t)+\frac{t}{|z|}\right)z.
\]
At \(t=0\), this is the identity.  At \(t=1\), it sends \(z\) to \(z/|z|\).  The missing origin is exactly what makes this formula possible without crossing \(0\).

Thus
\[
  \pi_1(\mathbb C\setminus\{0\})\cong\pi_1(S^1)\cong\mathbb Z.
\]
The hole is measured by how many times a loop winds around it.

\section{Functoriality: maps of spaces give homomorphisms}

A continuous based map
\[
  f:(X,x_0)\to(Y,y_0)
\]
sends a loop \(\gamma\) in \(X\) to the loop \(f\circ\gamma\) in \(Y\).  This respects homotopy and concatenation, so it induces a group homomorphism
\[
  f_*:\pi_1(X,x_0)\to\pi_1(Y,y_0).
\]
This is the real reason group homomorphisms appear in the same note as homotopy.  Algebraic topology assigns algebraic invariants to spaces and homomorphisms to continuous maps.

\section{Van Kampen as a computation machine}

The Seifert--van Kampen theorem says, roughly, that if a space is built from open pieces \(U\) and \(V\) with path-connected overlap, then \(\pi_1(U\cup V)\) is assembled from \(\pi_1(U)\) and \(\pi_1(V)\), with relations coming from the overlap.

For the figure-eight space \(S^1\vee S^1\), the result is
\[
  \pi_1(S^1\vee S^1)\cong F(a,b),
\]
the free group on two generators.  The two loops are independent because the wedge point imposes no commutation relation.

\section{Group completion and why inverses matter}

A monoid records composable operations with an identity; a group additionally allows reversal.  Loops are naturally reversible by traversing them backward, so the fundamental group is a group rather than merely a monoid.  In other contexts, one starts with a monoid and formally adds inverses; \(\mathbb N\) under addition completes to \(\mathbb Z\).

This mirrors the topological story: algebraic structures often arise by imposing equivalence relations and adding the formal operations needed to make the intended behavior exact.

\section{Loops, coverings, and algebraic invariants}

The central bridge is:
\[
  \text{continuous deformation of loops}
  \quad\longrightarrow\quad
  \text{homotopy classes}
  \quad\longrightarrow\quad
  \text{group elements}.
\]
The punctured plane shows the bridge concretely: a missing point creates winding number, winding number is an integer, and concatenation of loops becomes addition of integers.
